This study tested whether routinely available monitoring records could provide provisional monthly estimates of aquifer-system compaction while Tuku MLCW observations were delayed for six months. Hydraulic head changes, cGNSS vertical displacement, and local sediment composition were related to compaction in six standardized depth sections. A temporally ordered evaluation kept each six-month block of MLCW responses hidden until the block ended, preventing unavailable or future compaction observations from influencing the estimates.

The updated model achieved \placeholder{MAIN PERFORMANCE FINDING} and performed most clearly in \placeholder{BEST DEPTH SECTION OR SECTIONS}. Performance in \placeholder{WEAKEST DEPTH SECTION OR SECTIONS} remained limited by \placeholder{SUPPORTED OBSERVATIONAL LIMIT}. Periodic updating \placeholder{UPDATE FINDING}, while the error-calibrated intervals \placeholder{UNCERTAINTY FINDING}. These results \placeholder{FINAL CONCLUSION ABOUT BRIDGING DELAYED DATA DELIVERY}. The MLCW remains the reference measurement for evaluation and model updating. A schedule with only one field measurement every six or twelve months was evaluated separately as a no-update reduced-observation sensitivity analysis, in which the mean absolute endpoint error was 4.96~mm at a six-month horizon and 8.18~mm at a twelve-month horizon, and the mean signed endpoint error (bias) was 0.47~mm and 1.13~mm, respectively (\Cref{subsec:discussion_reduced_sampling}). Operational use of such a schedule would still require an update method that constrains the monthly estimates to the measured interval compaction.
